\subsubsection{Suggesting Applicability Conditions}
\label{app:condition-suggestions}
If TEPL cannot verify a rule under its declared condition, the compiler
can suggest an additional condition $G(p)$, rerun synthesis, and verify
the resulting program under the strengthened guard:
\begin{equation}
 \forall p,\quad \operatorname{Legal}_L(p)\land C(p)\land G(p)
 \Rightarrow p\in\operatorname{dom}(f)\land\operatorname{Valid}(p,f(p)).
 \label{eq:suggested-condition}
\end{equation}
TEPL checks that the suggested condition permits at least one legal
match and guarantees correctness. The search may fail to find a
condition or may find one stricter than necessary. The user must add
the suggestion to \code{where} and check the rule again.

The implemented search generates equalities between captured axis lists
and producer batch or free-group positions, together with identity
permutation restrictions. It tries single predicates and pairs from a
bounded candidate set. Each trial reruns construction and validation in
a fresh context under $C\land G$. A separate bounded concrete search
must find legal source metadata satisfying that condition. The witness
establishes nonvacuity; it is not evidence for universal correctness.
Failed witness search yields no suggestion. Deleting a redundant predicate
also requires fresh validation and a witness. The procedure claims neither
a weakest condition nor completeness of condition discovery.

For example, consider reassociation with omitted RHS attributes and no
\code{where} block. Imports are omitted:
\begin{lstlisting}[language=TEPL]
rule reassociate { (source: dot (inner: dot A B) C) => (dot A (dot B C)) }
\end{lstlisting}
Algebraic associativity alone is insufficient for generalized dots:
axis routing and output order also matter. Under an arithmetic policy
permitting reassociation, a sufficient restriction for this example is
\begin{lstlisting}[language=TEPL]
where {
  source.attrs.lhs_batching_dimensions == inner.axes.B.positions
}
\end{lstlisting}
The condition requires the outer dot's left batch positions to match
the positions of the inner dot's batch axes in its output, in the same
order. Here \code{inner.axes.B} denotes the inner operation's batch
group, not input tensor \code{B}. Appendix~\ref{extended:sec:partitioning} gives the batch-axis mappings.

After the user adds this sufficient condition, TEPL reconstructs and
validates the rule. Other valid matches may remain excluded.
